\documentclass[10pt,a4paper,fontset=fandol]{ctexart}
\usepackage[margin=2cm]{geometry}
\usepackage{needspace}
\ctexset{section={format=\large\sffamily\bfseries,beforeskip=14pt,afterskip=7pt},subsection={format=\normalsize\sffamily\bfseries,beforeskip=10pt,afterskip=5pt}}
\usepackage{fontspec,booktabs,longtable,array,xcolor,hyperref}
\setmainfont{TeX Gyre Pagella}
\setsansfont{TeX Gyre Heros}
\definecolor{reportblue}{HTML}{244D65}
\hypersetup{colorlinks=true,urlcolor=reportblue,linkcolor=reportblue}
\setlength{\parindent}{0pt}
\setlength{\parskip}{5pt}
\renewcommand{\arraystretch}{1.18}
\setlength{\tabcolsep}{5pt}
\emergencystretch=2em
\urlstyle{same}
\pagestyle{plain}
\title{\sffamily Jev 在 MoCa 因果与道德判断中的人类一致性}
\author{}
\date{2026 年 9 月 30 日}
\begin{document}
\begin{center}{\sffamily\LARGE Jev 在 MoCa 因果与道德判断中的人类一致性\par}\vspace{8pt}{\small 2026 年 9 月 30 日}\end{center}\vspace{6pt}
\section*{摘要}
在 MoCa 的 206 个英语故事上，Jev 与每题 25 人聚合判断的三类一致率为 43.69\%。因果题为 39.58\%，道德题为 53.23\%。Jev 能区分人类明确肯定与否定的题目，但对人类分歧较大的因果故事较少落入模糊区间。
\section{数据集}
材料包括 144 个因果故事与 62 个道德许可故事。参照为数据集每题的 25 个人类二元判断。三类为 Yes、No 和 Ambiguous，沿用 0.4／0.6 边界；Jev 的第三类由二选一概率派生。模型版本为 typesafe/jev-1.13-20260917。

因素名称沿用来源中的专家标注，同一故事可以包含多个因素。人类比例与 Jev 结构化选择概率具有不同含义，本文报告它们在这组故事上的一致性与数值差异。

\section{结果}
\Needspace{8\baselineskip}
\subsection{整体结果}
以数据集的人类聚合标签为一致性判据，不将人类多数判断称为道德正确答案。三类均匀随机参照为 33.33\%，总体多数类参照为 36.41\%。AUROC 只使用人类非模糊题。

{\small
\begin{longtable}{@{}>{\raggedright\arraybackslash}p{\dimexpr 0.2400\textwidth-2\tabcolsep\relax}>{\raggedright\arraybackslash}p{\dimexpr 0.1520\textwidth-2\tabcolsep\relax}>{\raggedright\arraybackslash}p{\dimexpr 0.1520\textwidth-2\tabcolsep\relax}>{\raggedright\arraybackslash}p{\dimexpr 0.1520\textwidth-2\tabcolsep\relax}>{\raggedright\arraybackslash}p{\dimexpr 0.1520\textwidth-2\tabcolsep\relax}>{\raggedright\arraybackslash}p{\dimexpr 0.1520\textwidth-2\tabcolsep\relax}@{}}
\toprule
\textbf{范围} & \textbf{题数} & \textbf{一致率} & \textbf{95\% 区间} & \textbf{AUROC} & \textbf{概率 MAE}\\
\midrule
\endfirsthead
\toprule
\textbf{范围} & \textbf{题数} & \textbf{一致率} & \textbf{95\% 区间} & \textbf{AUROC} & \textbf{概率 MAE}\\
\midrule
\endhead
全部 & 206 & 43.69\% & 36.89\%–50.49\% & 0.704 & 0.308\\
因果 & 144 & 39.58\% & 31.94\%–47.92\% & 0.748 & 0.333\\
道德 & 62 & 53.23\% & 40.32\%–66.13\% & 0.852 & 0.249\\
\bottomrule
\end{longtable}
}
\Needspace{8\baselineskip}
\subsection{判断分布}
{\small
\begin{longtable}{@{}>{\raggedright\arraybackslash}p{\dimexpr 0.2400\textwidth-2\tabcolsep\relax}>{\raggedright\arraybackslash}p{\dimexpr 0.1900\textwidth-2\tabcolsep\relax}>{\raggedright\arraybackslash}p{\dimexpr 0.1900\textwidth-2\tabcolsep\relax}>{\raggedright\arraybackslash}p{\dimexpr 0.1900\textwidth-2\tabcolsep\relax}>{\raggedright\arraybackslash}p{\dimexpr 0.1900\textwidth-2\tabcolsep\relax}@{}}
\toprule
\textbf{范围} & \textbf{来源} & \textbf{Yes} & \textbf{No} & \textbf{Ambiguous}\\
\midrule
\endfirsthead
\toprule
\textbf{范围} & \textbf{来源} & \textbf{Yes} & \textbf{No} & \textbf{Ambiguous}\\
\midrule
\endhead
因果 & 人类 & 48 & 50 & 46\\
因果 & Jev & 98 & 35 & 11\\
道德 & 人类 & 23 & 10 & 29\\
道德 & Jev & 22 & 21 & 19\\
\bottomrule
\end{longtable}
}
46 个被人类归为模糊的因果题中，Jev 只有 1 题同样归为模糊；道德题相应为 12/29。因果题中 Jev 的肯定判断明显更多，说明三类一致率较低主要涉及判断分布及边界差异。

\Needspace{8\baselineskip}
\subsection{因素分组的平均差异}
{\small
\begin{longtable}{@{}>{\raggedright\arraybackslash}p{\dimexpr 0.3200\textwidth-2\tabcolsep\relax}>{\raggedright\arraybackslash}p{\dimexpr 0.2267\textwidth-2\tabcolsep\relax}>{\raggedright\arraybackslash}p{\dimexpr 0.2267\textwidth-2\tabcolsep\relax}>{\raggedright\arraybackslash}p{\dimexpr 0.2267\textwidth-2\tabcolsep\relax}@{}}
\toprule
\textbf{分组变化} & \textbf{两组题数} & \textbf{人类差值} & \textbf{Jev 差值}\\
\midrule
\endfirsthead
\toprule
\textbf{分组变化} & \textbf{两组题数} & \textbf{人类差值} & \textbf{Jev 差值}\\
\midrule
\endhead
行动 → 不作为 & 80 / 32 & +0.028 & -0.179\\
知情 → 不知情 & 23 / 20 & -0.009 & +0.195\\
合取 → 析取 & 66 / 32 & -0.030 & -0.250\\
异常 → 正常 & 44 / 45 & -0.228 & -0.294\\
他人受益 → 自己受益 & 26 / 22 & +0.101 & -0.015\\
意外 → 工具性 & 18 / 30 & -0.118 & -0.067\\
可避免 → 不可避免 & 22 / 26 & +0.131 & +0.136\\
非个人施力 → 个人施力 & 24 / 24 & -0.058 & -0.067\\
\bottomrule
\end{longtable}
}
差值描述后一组相对前一组的平均肯定比例／概率变化。因素相互重叠，故事内容也随组变化，因此这些结果为描述性关联，不是单个因素的独立因果效应。

\Needspace{8\baselineskip}
\subsection{调用成本}
{\small
\begin{longtable}{@{}>{\raggedright\arraybackslash}p{\dimexpr 0.2400\textwidth-2\tabcolsep\relax}>{\raggedright\arraybackslash}p{\dimexpr 0.3800\textwidth-2\tabcolsep\relax}>{\raggedright\arraybackslash}p{\dimexpr 0.3800\textwidth-2\tabcolsep\relax}@{}}
\toprule
\textbf{输入 token} & \textbf{输出 token} & \textbf{已报告费用（美元）}\\
\midrule
\endfirsthead
\toprule
\textbf{输入 token} & \textbf{输出 token} & \textbf{已报告费用（美元）}\\
\midrule
\endhead
109,940 & 6,798 & 0.00461748\\
\bottomrule
\end{longtable}
}
206 次请求全部一次成功，没有未知费用或重试。

\Needspace{8\baselineskip}
\section{结论}
Jev 与人类聚合判断呈现部分相似性，道德题的一致性高于因果题。对人类明确肯定与否定的区分能力，并未转化为对人类模糊判断的同等识别。因素分组的部分平均方向相同，另一些方向不同；这些结果描述相同故事上的判断关系。
\Needspace{5\baselineskip}
\section*{参考文献}
\begin{enumerate}
\item Nie et al. (2023). MoCa: Measuring Human-Language Model Alignment on Causal and Moral Judgment Tasks. NeurIPS. \url{https://arxiv.org/abs/2310.19677}
\item MoCa official repository, commit 1b61a20294247480d64675ceb19751ef4e1e878f. \url{https://github.com/cicl-stanford/moca}
\end{enumerate}
\end{document}
